\begin{frame}{Aura Topic}
    
\end{frame}

\begin{frame}{Convenzioni elettromagnetiche: idea generale}
    \begin{columns}[T]
        \begin{column}{0.5\textwidth}
\small
\[
\boxed{\text{Le differenze stanno soprattutto nei fattori }4\pi,\ \varepsilon_0,\ \mu_0,\ c}
\]

\begin{itemize}
  \item \textbf{SI}: mantiene espliciti \(\varepsilon_0,\mu_0\).
  \item \textbf{cgs gaussiano}: assorbe \(\varepsilon_0,\mu_0\); compaiono fattori \(4\pi\).
  \item \textbf{Heaviside--Lorentz}: unità razionalizzate; i \(4\pi\) spariscono da Maxwell.
  \item \textbf{Unità naturali}: spesso \(\hbar=c=1\); usate in QFT/fisica nucleare.
\end{itemize}
        \end{column}
        \begin{column}{0.5\textwidth}
\small
\[
\begin{array}{c|c|c}
\text{sistema} & V(r) & \alpha \\ \hline
\text{SI} &
\dfrac{q_1q_2}{4\pi\varepsilon_0 r} &
\dfrac{e^2}{4\pi\varepsilon_0\hbar c}
\\[3mm]
\text{cgs gaussiano} &
\dfrac{q_1q_2}{r} &
\dfrac{e^2}{\hbar c}
\\[3mm]
\text{Heaviside--Lorentz} &
\dfrac{q_1q_2}{4\pi r} &
\dfrac{e^2}{4\pi\hbar c}
\end{array}
\]

In unità naturali \(\hbar=c=1\):
\[
\boxed{
V(r)=\frac{\alpha}{r}
}
\]

In fisica nucleare:
\[
\boxed{
\alpha\hbar c \simeq 1.44\ \mathrm{MeV\,fm}
}
\]
            
        \end{column}
    \end{columns}
    
\end{frame}

\begin{frame}{Maxwell: SI vs cgs gaussiano - Energia elettromagnetica}
    \begin{columns}[T]
        \begin{column}{0.5\textwidth}
\small
\[
\begin{array}{c|c|c}
 & \text{SI} & \text{cgs gaussiano} \\ \hline
\nabla\cdot\mathbf E &
\dfrac{\rho}{\varepsilon_0} &
4\pi\rho
\\[3mm]
\nabla\cdot\mathbf B &
0 & 0
\\[2mm]
\nabla\times\mathbf E &
-\dfrac{\partial\mathbf B}{\partial t} &
-\dfrac1c\dfrac{\partial\mathbf B}{\partial t}
\\[3mm]
\nabla\times\mathbf B &
\mu_0\mathbf J+\mu_0\varepsilon_0
\dfrac{\partial\mathbf E}{\partial t}
&
\dfrac{4\pi}{c}\mathbf J+
\dfrac1c\dfrac{\partial\mathbf E}{\partial t}
\end{array}
\]
        \end{column}
        \begin{column}{0.5\textwidth}
\small
\[
\begin{array}{c|c|c}
\text{sistema} & u_{\rm EM} & \mathbf S \\ \hline
\text{SI} &
\dfrac12\left(\varepsilon_0E^2+\dfrac{B^2}{\mu_0}\right)
&
\dfrac1{\mu_0}\mathbf E\times\mathbf B
\\[4mm]
\text{cgs gaussiano} &
\dfrac{E^2+B^2}{8\pi}
&
\dfrac{c}{4\pi}\mathbf E\times\mathbf B
\\[4mm]
\text{Heaviside--Lorentz} &
\dfrac12(E^2+B^2)
&
\mathbf E\times\mathbf B
\end{array}
\]

Per una sfera uniformemente carica in cgs:
\[
\boxed{
U_{\rm Coul}=\frac35\frac{Q^2}{R}
}
\]
            
        \end{column}
    \end{columns}
    
\end{frame}

\begin{frame}{Promemoria per fisica nucleare}
\small
\begin{itemize}
  \item In cgs/nucleare:
  \[
  e^2=\alpha\hbar c\simeq1.44\ \mathrm{MeV\,fm}.
  \]

  \item Energia Coulombiana di una sfera con \(Z\) protoni:
  \[
  \boxed{
  U_{\rm Coul}
  =
  \frac35\frac{Z^2e^2}{R}
  }
  \]

  \item In unità naturali:
  \[
  \hbar=c=1,
  \qquad
  e^2=
  \begin{cases}
  \alpha & \text{gaussiano},\\
  4\pi\alpha & \text{Heaviside--Lorentz}.
  \end{cases}
  \]

  \item Attenzione: quando un testo scrive \(e^2\), bisogna capire
  se usa convenzione gaussiana o razionalizzata.
\end{itemize}
\end{frame}

\begin{frame}{Fermi energy}
La definizione di energia di Fermi non \'e completamente universale.
    \begin{itemize}
        \item In fisica relativistica, nucleare e delle stelle compatte \'e comune definire l'energia di Fermi come l'energia totale della particella alla superficie di Fermi: $E_F=\sqrt{p_F^2c^2+m^2c^4}$ oppure, in unità naturali $\hbar=c=1$ $E_F=\sqrt{k_F^2+m^2}$.
Con questa convenzione la massa a riposo \'e inclusa. Per un gas di Fermi libero a temperatura nulla vale inoltre $\mu(T=0)=E_F$.
\item In trattazioni non relativi stiche, invece, con energia di Fermi si indica spesso la sola energia cinetica $\epsilon_F=\frac{p_F^2}{2m}$.
    L'analogo relativistico della sola parte cinetica \'e $E_F^{\mathrm{kin}}=\sqrt{p_F^2c^2+m^2c^4}-mc^2$, così che $E_F^{\mathrm{tot}}=mc^2+E_F^{\mathrm{kin}}$. Quando si considera il grado di degenerazione, la scala energetica rilevante \'e la parte cinetica dell'energia di Fermi, $k_B T_F = E_F^{\mathrm{kin}}$ e il regime fortemente degenere \'e caratterizzato da $\frac{T}{T_F}=\frac{k_B T}{E_F^{\mathrm{kin}}}\ll 1$. Usare invece l'energia di Fermi comprensiva della massa a riposo renderebbe fuorviante il confronto, soprattutto per particelle massive come protoni e neutroni.
    \end{itemize}
\end{frame}

\begin{frame}{Energia di un nucleo legato}

Per un sistema di \(A\) nucleoni, schematicamente: $E_{tot}=A m_N + T + V$ dove $T>0$ \'e l'energia cinetica interna dei nucleoni e \(V\) contiene le interazioni nucleari. Definiamo l'energia di legame rispetto a \(A\) nucleoni liberi a riposo: $E_b\equiv E_{tot}-A m_N=T+V$. Per un sistema legato: $E_b<0$; quindi l'energia cinetica dei nucleoni legati non è trascurata: il legame richiede semplicemente che $T+V<0$ ossia che il contributo attrattivo dell'interazione prevalga sull'energia cinetica. Con la convenzione alternativa \(B>0\) per la binding energy: $B=-E_b$.

Energia per nucleone e densit\'a di energia. Per materia nucleare uniforme a densit\'a barionica \(n_B\), $\varepsilon=\varepsilon_{rest}+\varepsilon_{kin}+\varepsilon_{int}$ \'e la densit\'a di energia totale; poich\'e $\varepsilon_{rest}=n_B m_N$, si ha $\frac{\varepsilon}{n_B}=m_N+\frac{T}{A}+\frac{V}{A}$ e quindi, per definizione, $\frac{E_b}{A}=\frac{\varepsilon}{n_B}-m_N=\frac{T}{A}+\frac{V}{A}$. 
Alla saturazione della symmetric nuclear matter: $\frac{E_b}{A}(n_0)\simeq -16~\mathrm{MeV}$ per cui $\frac{\varepsilon(n_0)}{n_0}\simeq m_N-16~\mathrm{MeV}$

Moto di Fermi ed energia di legame. Alla densit\'a di saturazione, $k_{F,0}\simeq 260~\mathrm{MeV}$. Nel semplice modello di gas di Fermi non relativistico: $T_F=\frac{k_F^2}{2m_N}\simeq 36~\mathrm{MeV}$, mentre l'energia cinetica media \'e $\frac{T}{A}=\frac{3}{5}T_F\simeq 22~\mathrm{MeV}$, quindi i nucleoni nel ground state hanno una notevole Fermi motion.

Poich\'e $\frac{E_b}{A}=\frac{T}{A}+\frac{V}{A}\simeq -16~\mathrm{MeV}$, schematicamente si pu\'o avere $\frac{T}{A}\sim +22~\mathrm{MeV}$, $\frac{V}{A}\sim -38~\mathrm{MeV}$. $E_b/A=-16~\mathrm{MeV}$ \'e il bilancio netto $T/A+V/A$ non il solo contributo potenziale.

\end{frame}

\begin{frame}{Constituents of NS: degenerate leptons, baryons and quarks.}
    \begin{itemize}
        \item Baade, Zwicky: Study of SN explosion when enough iron is made (endpoint exotermic reactions) the core collapse releasing gravitational binding energy: BE gravitazionale $10\% M_{NS}$, BE per nucleone per nuclei di Fe \SI{9}{\mega\ev}, $1\% M_{NS}$. 
        \item Luminous stars show Doppler broadening of spectral lines: collapsed core spun up (angular momentum conservation). Limits on rotation by mass loss at equator and gravitational radiation.
        \item Rotating relativistic stars set local inertial frames into rotation resp. to distant stars: infalling object from rst will acquire larger angular velocity as approaches.
        \item Dense many-body problems at level of nuclear physics, quarks and gluons: NS could be only object in which a coulomb lattice structure formed from two phases hadronic matter exists - WD are nuclei immersed in electron gas arranged in a Coulomb lattice.
        \item Strange stars - Strange matter hypothesis. Equilibrium energy per baryon of quarks up down strange is lower than the mass of nucleon, ie energy per nucleon of most bound nucleus (iron). Strange quark matter is the absolute ground state of strong interactions. Ipotesi di stabilit\'a richiede $(\frac{E}{A})_{uds,P=0}<\SI{930.4}{\mega\ev}$ valore riferimento per materia nucleare pi\'u stabile (energia libera di Gibbs per barione $g=\frac{\epsilon+P}{n_B}=\mu_B$ e per $P=0$ $g=\frac{E}{A}$, $\epsilon$ densit\'a energia totale).
    \end{itemize}
\end{frame}

\begin{frame}{Segnatura e segni in RG: $(-+++)\to(+---)$, $\tensor{g{'}}{^{\mu\nu}}\to-\tensor{g}{^{\mu\nu}}$}
    \begin{itemize}
        \item $\epsilon=$\lbt{+1 (+---)}{-1 (-+++)}: $ds^2=\epsilon c^2\,d\tau^2$, $\tensor{u}{^\mu}\tensor{u}{_\mu}=\epsilon c^2$.
        \item Intervallo timelike: $(+---): ds^2>0$, $(-+++): ds^2<0$
        \item Tempo proprio: $(+---): ds^2=c^2d\tau^2$, $(-+++): ds^2=-c^2d\tau^2$.
        \item Quadrivelocit\'a: $(+---): \tensor{u}{^\mu}\tensor{u}{_\mu}=c^2$, $(-+++): \tensor{u}{^\mu}\tensor{u}{_\mu}=-c^2$.
        \item Vettore spaziale: $\tensor{v}{^\mu}\tensor{v}{_\mu}<0$, $\tensor{v}{^\mu}\tensor{v}{_\mu}>0$.
        \item $\Gamma=\frac{1}{2}\tensor{g}{^{\rho\sigma}}(\tensor{\partial}{_\mu}\tensor{g}{_{\sigma\nu}}+\tensor{\partial}{_\nu}\tensor{g}{_{\sigma\mu}}-\tensor{\partial}{_\sigma}\tensor{g}{_{\mu\nu}})$, $\nabla$ non cambiano, $\tensor{V}{_\mu}$ cambia segno
        \item Invariato: $\tensor{R}{^\rho_{\sigma\mu\nu}}$, $\tensor{R}{_{\mu\nu}}=\tensor{R}{^\alpha_{\mu\alpha\nu}}$, $\tensor{G}{^{\mu\nu}}=\tensor{R}{^{\mu\nu}}-\frac{1}{2}\tensor{g}{^{\mu\nu}}R$
        \item Cambia Segno: $R=\tensor{g}{^{\mu\nu}}\tensor{R}{_{\mu\nu}}$
        \item Operatore D'alambert: $\square\phi=\tensor{g}{^{\mu\nu}}\tensor{\nabla}{_\mu}\tensor{\nabla}{_\nu}\phi$. Cambia segno: $\square'\phi=-\square\phi$.
        \item Tensore energia-impulso: le grandezze fisiche misurate non cambiano $\tensor{T}{_{00}}=\rho$, $\tensor{T}{_{ii}}=P$, in sistema ortonormale comovente. $\tensor{u}{^\mu}\tensor{u}{_\mu}=1$: $(+---)$ $\tensor{T}{_{\mu\nu}}=(\rho+P)\tensor{u}{_\mu}\tensor{u}{_\nu}-P\tensor{g}{_{\mu\nu}}$, $(-+++)$ $\tensor{T}{_{\mu\nu}}=(\rho+P)\tensor{u}{_\mu}\tensor{u}{_\nu}+P\tensor{g}{_{\mu\nu}}$
    \end{itemize}
\end{frame}

\begin{frame}{Solution of Einstein Eq.}
    \begin{itemize}
        \item Sch metric outside a static spherical star.
        \item Kerr metric outside rotating BH
        \item Einstein equation for interior solution of a spherically static star: Oppenheimer-Volkoff eqs. (numerically solved)

    \end{itemize}
\end{frame}

\begin{frame}{EOS e materia fortemente degenere}\frameinlbftrue
    \begin{block}{EOS: Relation between Pressure and Energy Density}
        \begin{itemize}
            \item Local properties: shear, bulk viscosity, dissipation, emissivity
            \item Global properties of bulk matter
        \end{itemize}
    \end{block}
    \begin{block}{Materia degenere fredda: $\frac{T}{T_{F,i}}\ll1$}
        \begin{align*}
            &K_BT_{F,i}=\sqrt{k_{F,i}^1+m_i^2}-m_i\\
            &T=0\tag{approx}
        \end{align*}
    \end{block}
\end{frame}
